%% ----------------------------------------------------------------------------
%%  §5.1  Tasso di accrescimento. Significato meccanico della derivata
%% ----------------------------------------------------------------------------
\section{Tasso di accrescimento. Significato meccanico della derivata}
\label{sec:tasso}

Consideriamo un semplice processo di crescita di un corpo, supponendo che il peso $p$
del corpo vari al crescere del tempo; cioè supponiamo che il peso sia una funzione del
tempo, $p=p(t)$. Prendiamo in considerazione le variazioni di peso a partire da un certo
istante $t$. Al tempo $t$ il peso è $p(t)$, mentre all'istante $t+h$, dopo che è
trascorso un tempo uguale ad $h$, il peso è $p(t+h)$. Quindi, nell'intervallo di tempo
$h$, il cambio di peso è $p(t+h)-p(t)$. Il rapporto
\begin{equation}\label{eq:tasso-medio}
  \frac{p(t+h)-p(t)}{h}
\end{equation}
dà un'indicazione di quanto sia cambiato il peso per unità di tempo. Più precisamente
il rapporto \eqref{eq:tasso-medio} esprime la variazione \emph{media} per unità di tempo
del peso nell'intervallo $[t,t+h]$ (o $[t+h,t]$ se $h<0$). Si chiama anche \emph{tasso
medio di accrescimento}, o tasso medio di variazione, o \emph{velocità media di
accrescimento}.

Invece della variazione media nell'intervallo $[t,t+h]$, spesso è più utile la
variazione \emph{istantanea} al tempo $t$. Intuitivamente si considera l'espressione
\eqref{eq:tasso-medio} per alcuni valori di $h$ sempre più vicini a zero. Con il
linguaggio più preciso introdotto nei capitoli precedenti, si calcola il limite, per
$h\to0$, del rapporto \eqref{eq:tasso-medio}:
\begin{equation}\label{eq:tasso}
  \text{tasso di accrescimento}=\lim_{h\to0}\frac{p(t+h)-p(t)}{h}.
\end{equation}
Si noti come sia indispensabile usare il limite per $h\to0$, dato che, se nella
\eqref{eq:tasso-medio} si pone direttamente $h=0$, si ottiene un'espressione senza
significato. Invece, con il linguaggio dei limiti, il limite in \eqref{eq:tasso} è una
forma indeterminata $\frac00$ (se $p$ è una funzione continua), del tipo già incontrato
nell'esempio della velocità istantanea del \cref{sec:premessa-funzioni}.

Per mostrare come il tasso di accrescimento possa essere effettivamente calcolato,
consideriamo a titolo di esercizio il caso in cui il peso dipenda dal tempo in modo
quadratico,
\begin{equation}\label{eq:peso-quadratico}
  p(t)=t^{2};
\end{equation}
per ogni $h\neq0$ risulta
\begin{equation}\label{eq:tasso-t2}
  \frac{p(t+h)-p(t)}{h}=\frac{(t+h)^{2}-t^{2}}{h}=\frac{t^{2}+2ht+h^{2}-t^{2}}{h}=2t+h .
\end{equation}
L'ultimo membro della \eqref{eq:tasso-t2} tende a $2t$ per $h\to0$. Quindi il tasso di
accrescimento vale $2t$. Ciò significa che, al crescere del tempo $t>0$, non soltanto il
peso cresce come $t^{2}$, ma anche il \emph{cambiamento} di peso per unità di tempo
aumenta (nel caso in considerazione, in modo proporzionale al tempo).

Proponiamo un esempio numerico: secondo la legge $p(t)=t^{2}$, al tempo $t=10$ il peso
risulta essere uguale a $p(10)=100$. Il tasso di accrescimento, uguale a $2t$, al tempo
$t=10$ vale $20$. Ciò significa che, dopo una unità di tempo, il peso del corpo aumenta
di circa $20$ unità; quindi $p(11)$ vale all'incirca $120$. Si noti che effettivamente il
valore trovato, $120$, non differisce di molto da $p(11)=11^{2}=121$; approfondiremo
questo aspetto nel \cref{sec:retta-tangente} e, in modo più completo, con la formula di
Taylor. Il programma seguente calcola i tassi medi per $h$ sempre più piccoli, in
aritmetica razionale esatta.

\begin{python}[Tassi medi di accrescimento per $p(t)=t^{2}$ e $t=10$]
from fractions import Fraction as F

p = lambda t: t**2                    # peso in funzione del tempo
t = F(10)
for k in range(0, 7):
    h = F(1, 10**k)                   # aritmetica razionale esatta
    tasso_medio = (p(t + h) - p(t)) / h
    print(f"h = {str(h):>9}   tasso medio = {float(tasso_medio)}")
print("previsione lineare p(11) ~", p(t) + 2 * t * 1,
      "  valore vero:", p(t + 1))
\end{python}
\begin{risultato}
h =         1   tasso medio = 21.0
h =      1/10   tasso medio = 20.1
h =     1/100   tasso medio = 20.01
h =    1/1000   tasso medio = 20.001
h =   1/10000   tasso medio = 20.0001
h =  1/100000   tasso medio = 20.00001
h = 1/1000000   tasso medio = 20.000001
previsione lineare p(11) ~ 120   valore vero: 121
\end{risultato}

I tassi medi valgono esattamente $20+h$, come previsto dalla \eqref{eq:tasso-t2}, e si
avvicinano a $20$; ma, ancora una volta, il calcolo con alcuni valori di $h$ suggerisce
il limite, mentre è la \eqref{eq:tasso-t2} a dimostrarlo.

\subsection{La velocità istantanea}

Abbiamo già detto che «velocità di accrescimento» è sinonimo di «tasso di
accrescimento»; ciò deriva dal fatto che una velocità si definisce in modo analogo a
quanto fatto sopra. Consideriamo ad esempio un velivolo che percorre una rotta
rettilinea, e indichiamo con $s(t)$ lo spazio percorso in funzione del tempo. La
\emph{velocità media} nell'intervallo di tempo $[t,t+h]$ è uguale al rapporto tra lo
spazio percorso $s(t+h)-s(t)$ e il tempo $h$ impiegato a percorrerlo. La \emph{velocità
istantanea} (quella che, a meno del vento, si ricava dagli strumenti di bordo) è il
limite, per $h\to0$, della velocità media; quindi
\begin{equation}\label{eq:velocita-istantanea}
  \text{velocità istantanea}=\lim_{h\to0}\frac{s(t+h)-s(t)}{h}.
\end{equation}
Il significato geometrico è illustrato nella \cref{fig:velocita-media}: la velocità media
è la pendenza della retta che congiunge due punti del grafico della legge oraria; la
velocità istantanea è la pendenza limite di tali rette quando il secondo punto si
avvicina al primo.

\begin{figure}[htb]
\centering
\begin{tikzpicture}
\begin{axis}[width=0.9\textwidth, height=5.6cm, xmin=0, xmax=6, ymin=0, ymax=5.6,
  axis lines=left, xtick={2,4.6}, xticklabels={$t$,$t+h$}, ytick={1.404,4.6221},
  yticklabels={$s(t)$,$s(t+h)$}, xlabel={tempo}, ylabel={spazio}, clip=false]
  \addplot[curva, domain=0.2:5.2, samples=100] {0.5+0.25*x^2-0.012*x^3};
  \addplot[tratteggio] coordinates {(0,1.404) (2,1.404) (2,0)};
  \addplot[tratteggio] coordinates {(0,4.6221) (4.6,4.6221) (4.6,0)};
  \addplot[curva2, domain=1.2:5.1] {1.404+(4.6221-1.404)/2.6*(x-2)};
  \addplot[curva3, domain=0.4:4.2] {1.404+0.856*(x-2)};
  \addplot[only marks, mark=*, mark size=1.6pt, blu] coordinates {(2,1.404) (4.6,4.6221)};
  \node[font=\scriptsize, rame, anchor=west] at (axis cs:5.1,5.0) {secante};
  \node[font=\scriptsize, verde!50!black, anchor=north west] at (axis cs:3.6,2.35) {tangente};
\end{axis}
\end{tikzpicture}
\caption{La velocità media nell'intervallo $[t,t+h]$ è la pendenza della retta secante
(rame) per i punti $(t,s(t))$ e $(t+h,s(t+h))$ del grafico della legge oraria; quando
$h\to0$ la secante tende alla retta tangente (verde), la cui pendenza è la velocità
istantanea.}
\label{fig:velocita-media}
\end{figure}

È chiaro che nei due esempi precedenti lo schema matematico è identico. In entrambi gli
esempi occorre calcolare \emph{il limite di un rapporto incrementale}, così chiamato
perché a denominatore c'è l'incremento $h$ della variabile indipendente, mentre a
numeratore c'è il corrispondente incremento della variabile dipendente. Occorre
calcolare il limite del rapporto incrementale anche in molte altre situazioni, analoghe
a quelle dei due esempi esposti:
\begin{itemize}
  \item l'\emph{accelerazione} di un corpo in moto rettilineo è il limite del rapporto
  incrementale della velocità rispetto al tempo;
  \item la \emph{portata in massa} di combustibile consumato da un motore è il limite del
  rapporto incrementale della massa di combustibile bruciata rispetto al tempo;
  \item la \emph{densità lineare} di una trave non omogenea è il limite del rapporto tra
  la massa di un tratto di trave e la sua lunghezza, quando la lunghezza tende a zero;
  \item il \emph{gradiente termico} verticale dell'atmosfera è il limite del rapporto
  incrementale della temperatura rispetto alla quota.
\end{itemize}
Un altro esempio, di tipo geometrico, è studiato nel \cref{sec:retta-tangente}.
Introdurremo nel prossimo paragrafo la \emph{derivata} come limite del rapporto
incrementale, quando l'incremento tende a zero.

\begin{idea}[Il variometro: uno strumento che misura una derivata]
Il variometro, presente sul cruscotto di ogni aliante e di quasi ogni velivolo, indica
la \emph{velocità verticale}, cioè il tasso di variazione della quota rispetto al tempo,
$\lim_{h\to0}\frac{z(t+h)-z(t)}{h}$. Il variometro meccanico realizza fisicamente
un rapporto incrementale: una capsula, collegata direttamente alla presa statica, è
racchiusa in un serbatoio che comunica con la stessa presa soltanto attraverso un
sottile tubo capillare, che rallenta l'equalizzazione delle pressioni. Quando la quota
cambia, la pressione nella capsula segue subito quella esterna, mentre quella nel
serbatoio resta indietro, e la differenza di pressione che si stabilisce è, con buona
approssimazione, proporzionale alla rapidità di variazione della pressione statica, e
quindi della quota. Lo strumento ha un ritardo, perché in effetti confronta valori in
istanti diversi: misura cioè un rapporto incrementale con un incremento di tempo
piccolo ma non nullo. I variometri elettronici calcolano il rapporto incrementale per
via numerica, con gli accorgimenti che vedremo nel \cref{sec:def-derivata}.
\end{idea}

\begin{notastorica}[Newton, Leibniz e le notazioni]
Il calcolo differenziale nacque due volte, in modo indipendente. Isaac Newton lo
sviluppò negli anni 1665--1666 come «metodo delle flussioni». Per lui le grandezze sono
quantità che «fluiscono» nel tempo, le \emph{fluenti}; la loro velocità di variazione è
la \emph{flussione}, e il prodotto della flussione per un intervallo di tempo
infinitamente piccolo $o$ è il \emph{momento}, cioè l'incremento infinitesimo della
fluente. Il problema centrale, nel trattato \emph{De methodis serierum et fluxionum}
scritto nel 1671, è proprio il nostro: data una relazione tra quantità fluenti, trovare
la relazione tra le loro flussioni. Poiché la variabile indipendente è sempre il tempo,
la derivata di $y$ rispetto a $x$ compare come rapporto tra le flussioni di $y$ e di $x$.
Il trattato fu pubblicato soltanto nel 1736, in traduzione inglese (\emph{The Method of
Fluxions and Infinite Series}); la notazione col punto, $\dot x$, ancora oggi d'uso in
meccanica per le derivate rispetto al tempo, Newton la adottò soltanto alla fine del 1691.
Gottfried Wilhelm Leibniz pubblicò per primo, nell'ottobre 1684, la memoria citata
nell'epigrafe, uscita a Lipsia negli \emph{Acta Eruditorum}: poche pagine che rendevano
pubblici i differenziali $dx$, $dy$, ideati nei manoscritti parigini del 1675, e le regole
per calcolarli. Leibniz pensava le grandezze come formate da parti infinitesime, e la
sua notazione, che conserva nel simbolo $\frac{dy}{dx}$ l'origine della derivata come
rapporto di incrementi, si è dimostrata straordinariamente efficace, come vedremo con la
regola della catena. Sulla disputa per la priorità che seguì si veda la nota storica
del \S\ref{sec:teorema-fondamentale}; il nome «derivata» e il simbolo $f'(x)$ sono più
tardi (si veda la nota storica alla fine del \S\ref{sec:def-derivata}).
\end{notastorica}
